% Part: first-order-logic
% Chapter: sequent-calculus
% Section: rules-and-proofs

\documentclass[../../../include/open-logic-section]{subfiles}

\begin{document}

\iftag{FOL}
      {\olfileid{fol}{seq}{rul}}
      {\olfileid{pl}{seq}{rul}}

\olsection{Aturan lan \usetoken{P}{derivation}}

Kanggo andharan sabanjure, anggepen $\Gamma, \Delta, \Pi, \Lambda$
makili urutan winates saka !!{sentence}s.

\begin{defn}[Sekuen]
\emph{Sekuen} yaiku ekspresi kanthi wujud
\[
\Gamma \Sequent \Delta
\]
ing ngendi $\Gamma$ lan $\Delta$ iku urutan winates (bisa uga kosong)
saka !!{sentence}s ing basa $\Lang L$. $\Gamma$ diarani
\emph{anteseden}, dene $\Delta$ diarani \emph{sukseden}.
\end{defn}

\begin{explain}
Gagasan intuitif ing walike sekuen yaiku: yen kabeh !!{sentence}s ing
anteseden bener, saora-orane salah siji saka !!{sentence}s ing sukseden
bener. Yaiku, yen $\Gamma = \tuple{!A_1, \dots, !A_m}$ lan
$\Delta = \tuple{!B_1, \dots, !B_n}$, mula $\Gamma \Sequent \Delta$
bener yen lan mung yen
\[
(!A_1 \land \cdots \land !A_m) \lif (!B_1 \lor \cdots \lor
!B_n)
\]
bener. Ana rong kasus mligi: nalika $\Gamma$ kosong lan nalika
$\Delta$~kosong. Nalika $\Gamma$ kosong, yaiku $m = 0$, $\quad
\Sequent \Delta$ bener yen lan mung yen $!B_1 \lor \dots \lor !B_n$
bener. Nalika $\Delta$ kosong, yaiku $n = 0$, $\Gamma \Sequent \quad$
bener yen lan mung yen $\lnot(!A_1 \land \dots \land !A_m)$ bener.
Kita ngarani sekuen valid yen lan mung yen !!{sentence} kang cocog
iku valid.
\end{explain}

Yen $\Gamma$ iku urutan !!{sentence}s, kita nulis $\Gamma, !A$ kanggo
asil nambahake $!A$ ing pucuk tengen~$\Gamma$ (lan $!A, \Gamma$
kanggo asil nambahake $!A$ ing pucuk kiwa~$\Gamma$). Yen $\Delta$ uga
urutan !!{sentence}s, mula $\Gamma, \Delta$ iku sambungan rong urutan
kasebut.

\begin{defn}[Sekuen Wiwitan]
\emph{Sekuen wiwitan} yaiku sekuen
\iftag{prvFalse,prvTrue}{kanthi salah siji wujud iki:
  \begin{enumerate}
    \item $!A \Sequent !A$
    \tagitem{prvTrue}{$\quad \Sequent \ltrue$}{}
    \tagitem{prvFalse}{$\lfalse \Sequent \quad$}{}
  \end{enumerate}}
{kanthi wujud $!A \Sequent !A$} kanggo saben !!{sentence} $!A$ ing basa kasebut.
\end{defn}

!!^{derivation}s ing kalkulus sekuen iku wit-wit sekuen tartamtu.
Sekuen paling dhuwur yaiku sekuen wiwitan, lan yen sawijining sekuen
mapan ing sangisore siji utawa rong sekuen liyane, sekuen iku kudu
ndherek kanthi bener miturut aturan inferensi. Aturan kanggo
$\Log{LK}$ dipérang dadi rong jinis utama: aturan \emph{logis} lan
aturan \emph{struktural}. Aturan logis dijenengi miturut !!{main
  operator} saka !!{sentence} kang ngemot $!A$ lan/utawa $!B$ ing
sekuen ngisor. Saben aturan duwe rong versi, siji kanggo nginferensi
sekuen kanthi !!{sentence} kang ngemot !!{operator} ing sisih kiwa,
lan siji maneh kanthi !!{sentence} kasebut ing sisih tengen.

\end{document}
